% problem-type: truyen-song-3d-cauchy-co-nguon
% order: 30
% Bai toan Cauchy phuong trinh truyen song 3D khong thuan nhat co nguon
% Bo cuc: De vi du -> Cong thuc -> De + Bai giai

\section{Truyền sóng 3D có nguồn --- đối xứng phẳng (hạ 1 biến)}

% ---------------------------------------------------------------
\subsection*{Đề bài ví dụ}
% ---------------------------------------------------------------
Xét bài toán Cauchy cho phương trình truyền sóng
\[
  u_{tt}(x,y,z,t) = \Delta u(x,y,z,t) + \chi_{[0,1]}(t)\,\chi_{[-1,1]}(x),
  \quad (x,y,z)\in\mathbb{R}^3,\ t>0,
\]
với điều kiện ban đầu $u(x,y,z,0)=0$ và
\[
  u_t(x,y,z,0) = \chi_{B_1(0,0,2)}(x,y,z) - \chi_{B_1(0,0,-2)}(x,y,z),
\]
trong đó $B_r(x,y,z)$ là hình cầu tâm $(x,y,z)$ bán kính $r$.
\textbf{Hãy tính $u(100,y,0,t)$ với $y\in\mathbb{R}$, $t>0$.}

\emph{\small(Nguồn: Đề thi MAT2306 --- PTĐHR 1, HK2 2024--2025, K67, Câu 2. Có đáp án chính thức.)}

% ---------------------------------------------------------------
\subsection*{Nhận ra dạng này khi nào?}
% ---------------------------------------------------------------
\begin{itemize}
  \item PDE dạng $u_{tt}-\Delta u = f(x,t)$ trên $\mathbb{R}^3\times(0,\infty)$, vế phải $f\neq0$ (có nguồn).
  \item Dữ liệu ban đầu cho trên toàn $\mathbb{R}^3$ (không có biên không gian) $\Rightarrow$ bài toán Cauchy.
  \item Số chiều $n=3$ (lẻ) $\Rightarrow$ dùng được nguyên lý Huygens mạnh.
\end{itemize}
\begin{meobox}
Hai mẹo cứu cánh: (1) nếu điểm hỏi nằm trên mặt phẳng \emph{đối xứng} của dữ liệu, phần
dữ liệu thường triệt tiêu; (2) nếu nguồn không phụ thuộc $y,z$ thì bài toán \emph{hạ về 1D}.
\end{meobox}
\begin{luuybox}
Đề thực thường \textbf{TRỘN} nhiều phần đối xứng khác nhau. Chồng chất tách $u=\sum u_j$, mỗi $u_j$
chọn đúng cách hạ chiều: phần phụ thuộc \emph{một biến phẳng} $\to$ mục này; phần \emph{đối xứng cầu}
($\chi_{B_\rho}$, theo $r$) $\to$ mục ``đối xứng cầu ($v=ru$)''.
\end{luuybox}

% ---------------------------------------------------------------
\subsection*{Công thức \& phương pháp}
% ---------------------------------------------------------------

\subsubsection*{Chồng chất: tách bài toán}
% method: nguyen-ly-chong-chat
\begin{dieukienbox}
Dùng khi: PDE \textbf{tuyến tính} --- tách bài toán thành nhiều phần đơn giản hơn,
giải từng phần rồi cộng lại.
\end{dieukienbox}

\begin{nhobox}
\textbf{Nguyên lý chồng chất}: với $u_{tt}-\Delta u = f$, $u(0)=g$, $u_t(0)=h$, tách theo VAI TRÒ:
\[
  u = \underbrace{u^{(g)} + u^{(h)}}_{\text{phần do DỮ LIỆU (nguồn tắt)}} + \underbrace{u^{(f)}}_{\text{phần do NGUỒN (dữ liệu tắt)}}.
\]
\begin{itemize}
  \item $u^{(g)}$: $u_{tt}=\Delta u$ (không nguồn), $u(0)=g$, $u_t(0)=0$.
  \item $u^{(h)}$: $u_{tt}=\Delta u$ (không nguồn), $u(0)=0$, $u_t(0)=h$.
  \item $u^{(f)}$: $u_{tt}-\Delta u=f$ (có nguồn), $u(0)=0$, $u_t(0)=0$ --- dùng Duhamel.
\end{itemize}
\textbf{Luôn kiểm tra}: cộng các bài con phải ra lại đúng PDE $+$ dữ liệu của bài gốc.
\end{nhobox}

\begin{meobox}
Nếu $h = h_1 - h_2$ (hiệu hai hàm), tách tiếp: $u^{(h)} = u^{(h_1)} - u^{(h_2)}$.
Ví dụ: $h = \chi_{B_1(0,0,2)} - \chi_{B_1(0,0,-2)}$ $\Rightarrow$ hai bài toán con, xử lý từng quả cầu.
\end{meobox}


\subsubsection*{Công thức Kirchhoff (phần dữ liệu ban đầu)}
% method: cong-thuc-kirchhoff
\begin{dieukienbox}
Dùng khi: phương trình sóng \textbf{thuần nhất} ($u_{tt}=c^2\Delta u$, không nguồn) trên $\mathbb{R}^3$.
Chỉ đúng cho $n=3$ chiều không gian. Vận tốc lan truyền là $c$.
\end{dieukienbox}

\begin{nhobox}
\textbf{Dữ liệu ban đầu:} $g(x)=u(x,0)$ (vị trí lúc đầu), $\;h(x)=u_t(x,0)$ (vận tốc lúc đầu).

\textbf{Công thức Kirchhoff} ($n=3$, vận tốc $c$) --- tích phân trên mặt cầu \textbf{bán kính $ct$}:
\[
  u(x,t)=\partial_t\!\left(\frac{1}{4\pi c^2 t}\int_{\partial B_{ct}(x)} g\,dS\right)
  +\frac{1}{4\pi c^2 t}\int_{\partial B_{ct}(x)} h\,dS .
\]
Mặt cầu $\partial B_{ct}(x)$ tâm $x$ bán kính $ct$ có diện tích $4\pi(ct)^2=4\pi c^2t^2$, nên
$\dfrac{1}{4\pi c^2t^2}\int_{\partial B_{ct}(x)}\!\cdots\,dS$ chính là \textbf{giá trị trung bình} trên mặt cầu đó.

Trường hợp $g=0$ (đứng yên lúc đầu) $\Rightarrow$ chỉ còn:
\[
  u(x,t)=\frac{1}{4\pi c^2 t}\int_{\partial B_{ct}(x)} h\,dS .
\]
\end{nhobox}

\begin{luuybox}
Đừng quên $c$: bán kính mặt cầu là $ct$ (\emph{không} phải $t$), hệ số là $\dfrac{1}{4\pi c^2t}$. Khi $c=1$ mới rút về $\dfrac{1}{4\pi t}$, mặt cầu bán kính $t$. (Đề thật hay gặp $c^2=4,9$ tức $c=2,3$.)
\end{luuybox}

\begin{meobox}
Nghiệm tại $(x,t)$ \emph{chỉ} phụ thuộc dữ liệu nằm \emph{đúng trên mặt cầu} $|y-x|=ct$ (không phải bên trong).
Vậy nếu mặt cầu này không chạm vùng mà $g,h$ khác $0$ thì $u(x,t)=0$.
\end{meobox}

\emph{(Cách tính tích phân mặt trên mặt cầu: xem Phụ lục C --- Kỹ thuật tính toán.)}

\subsubsection*{Nguyên lý Duhamel (phần nguồn)}
% method: nguyen-ly-duhamel
\begin{dieukienbox}
Dùng khi: phương trình sóng \textbf{có nguồn} $u_{tt}-c^2\Delta u=f$,
dữ liệu ban đầu bằng $0$ (hoặc đã tách phần do dữ liệu ban đầu bằng chồng chất).
\end{dieukienbox}

\begin{nhobox}
\textbf{Định nghĩa $w(x,t;s)$:} với mỗi $s$ cố định, $w$ giải phương trình \emph{không có nguồn}
$w_{tt}=c^2\Delta w$, nhưng đặt dữ liệu tại $t=s$ với \textbf{vận tốc đầu $=f(\cdot,s)$} (lát cắt nguồn):
\[
  \begin{cases}
    w_{tt} = c^2\Delta w & (t>s),\\
    w(x,s;\,s) = 0 & (\text{vị trí đầu}=0),\\
    w_t(x,s;\,s) = f(x,s) & (\text{vận tốc đầu}=f\neq0).
  \end{cases}
\]
\textbf{Nguyên lý Duhamel} --- gộp theo $s$:
\[
  u(x,t) = \int_0^t w(x,t;\,s)\,ds .
\]
\end{nhobox}

\begin{nhobox}
\textbf{Tìm $w$ thế nào:} dùng đúng công thức nghiệm \emph{không nguồn} đã biết (Kirchhoff/d'Alembert),
với dữ liệu $g=0$, $h=f(\cdot,s)$, và thay thời gian $t\to t-s$ (vì $w$ chỉ phụ thuộc thời gian trôi $t-s$).
Với vận tốc lan truyền $c$ (phương trình $u_{tt}=c^2\Delta u+f$):
\[
  \text{1D (d'Alembert): }\ w=\frac{1}{2c}\!\int_{x-c(t-s)}^{x+c(t-s)}\! f(\xi,s)\,d\xi,
\]
\[
  \text{3D (Kirchhoff): }\ w=\frac{1}{4\pi c^2(t-s)}\int_{\partial B_{c(t-s)}(x)}\! f(y,s)\,dS(y).
\]
(Hệ số vận tốc là $\tfrac{1}{2c}$, \emph{không} phải $\tfrac12$; khi $c=1$ thì hai cái trùng nhau.)
\end{nhobox}

\begin{luuybox}
\textbf{Đừng nhầm hai chữ ``không nguồn''!}
\begin{itemize}
  \item \emph{Phần do dữ liệu} $u^{(g)},u^{(h)}$ (chồng chất): không nguồn $+$ dữ liệu THẬT $(g,h)$. Có thể $=0$ ở vài điểm (vd do đối xứng) --- không liên quan Duhamel.
  \item \emph{Sóng con $w$} (Duhamel): không nguồn nhưng dữ liệu đầu $=f(\cdot,s)\neq0$, nên $w\neq0$. Đây mới là thứ Duhamel cộng lại.
\end{itemize}
``Không nguồn'' nói về \textbf{phương trình} (vế phải $=0$), KHÔNG phải dữ liệu $=0$.
\end{luuybox}

\begin{meobox}
$f$ có giá compact theo $s$ (vd $s\in[0,1]$) $\Rightarrow$ tích phân Duhamel tự thu hẹp về đoạn đó.
\end{meobox}


\subsubsection*{Hạ thấp chiều (nguồn độc lập $y,z$)}
% method: phuong-phap-ha-thap
\begin{dieukienbox}
Dùng khi: dữ liệu hoặc nguồn \textbf{không phụ thuộc một hoặc vài biến không gian}.
Ví dụ: nguồn $f(x,t)=\chi_{[-1,1]}(x_1)$ không phụ thuộc $y,z$
$\Rightarrow$ bài toán 3D thực ra chỉ ``sống'' trong 1D theo $x_1$.
\end{dieukienbox}

\begin{nhobox}
\textbf{Ý tưởng hạ thấp}: nếu mọi thứ chỉ phụ thuộc $x_1$ (không có $y,z$),
thì $\Delta u = u_{x_1 x_1}$ và bài toán trở thành phương trình sóng \textbf{1D}.

Áp dụng công thức d'Alembert cho bài 1D, hoặc suy nghiệm 2D/1D từ công thức 3D (Kirchhoff) bằng cách tích phân thêm theo biến phụ.
\end{nhobox}

\begin{meobox}
Dấu hiệu nhận ra: $f$ hoặc $h$ viết theo $\chi_{[-1,1]}(x_1)$ --- chỉ có $x_1$, không có $y,z$.
Khi đó có thể bỏ hai biến $y,z$ và giải bài 1D nhanh hơn nhiều.
\end{meobox}

\emph{(Sau khi hạ về 1D, dùng công thức d'Alembert --- xem dạng bài ``Truyền sóng 1D''. Nếu dữ liệu/nguồn
\textbf{đối xứng cầu} (chỉ phụ thuộc $r=|x|$) thì hạ về 1D bằng đổi biến $v=ru$ --- xem Phụ lục C.)}

\subsubsection*{Miền phụ thuộc / vận tốc hữu hạn}
% method: phan-tich-mien-phu-thuoc
\begin{dieukienbox}
Dùng khi: cần xác định nghiệm \emph{có bằng $0$ hay không} tại một điểm xa nguồn/dữ liệu.
Áp dụng cho mọi PDE hyperbolic vận tốc hữu hạn ($c=1$ ở đây).
\end{dieukienbox}

\begin{nhobox}
\textbf{Quy trình 3 bước:}
\begin{enumerate}
  \item Xác định giá (\emph{support}) của $g$, $h$, $f$.
  \item Tính khoảng cách từ điểm quan sát $x_0$ đến giá:
        $d_{\min} = \operatorname{dist}(x_0, \operatorname{supp})$.
  \item Kết luận: $u(x_0,t)=0$ khi $t < d_{\min}$ (sóng chưa tới);
        tín hiệu xuất hiện từ $t = d_{\min}$.
\end{enumerate}
Trong 3D (nguyên lý Huygens mạnh): tín hiệu cũng tắt sau $t = d_{\max}$
($d_{\max}$: khoảng cách lớn nhất đến giá).
\end{nhobox}

\begin{meobox}
Ví dụ: $x_0=(100,0,0)$, $\operatorname{supp}(h)\subset B_1(0,0,\pm 2)$.
Khoảng cách gần nhất $\approx 97$, xa nhất $\approx 103$.
$\Rightarrow$ $u^{(h)}(x_0,t)=0$ ngoài khoảng $t\in[97,103]$ (tính chính xác theo hình học).

Cho nguồn $f=\chi_{[0,1]}(s)\chi_{[-1,1]}(x_1)$: $\operatorname{dist}(x_0,\operatorname{supp}(f(\cdot,s)))\ge 99$,
cần $t-s\ge 99$ $\Rightarrow$ tín hiệu từ $f$ bắt đầu từ $t\approx 99$.
\end{meobox}


% Khái niệm dùng -> đã gom về Chương 0 (giả định đã biết).

% ---------------------------------------------------------------
\subsection*{Đề bài \& bài giải}
% ---------------------------------------------------------------

\paragraph{Nhắc lại đề.}
$u_{tt}=\Delta u+\chi_{[0,1]}(t)\chi_{[-1,1]}(x)$ trên $\mathbb{R}^3$, $u(\cdot,0)=0$,
$u_t(\cdot,0)=\chi_{B_1(0,0,2)}-\chi_{B_1(0,0,-2)}$. Tính $u(100,y,0,t)$.

\paragraph{Bước 0 --- Chồng chất (tách thành 2 bài con).} Tuyến tính nên bài gốc
($u_{tt}=\Delta u+f$, $u(\cdot,0)=0$, $u_t(\cdot,0)=h$) tách $u=u_h+u_f$:
\[
  (A)\ \begin{cases}(u_h)_{tt}=\Delta u_h\\ u_h(\cdot,0)=0,\ (u_h)_t(\cdot,0)=h\end{cases}
  \quad
  (B)\ \begin{cases}(u_f)_{tt}=\Delta u_f+f\\ u_f(\cdot,0)=0,\ (u_f)_t(\cdot,0)=0\end{cases}
\]
với $h=\chi_{B_1(0,0,2)}-\chi_{B_1(0,0,-2)}$, $f=\chi_{[0,1]}(t)\chi_{[-1,1]}(x)$.
Kiểm tra: cộng lại được đúng bài gốc ($u_h$ gánh dữ liệu đầu, $u_f$ gánh nguồn).
Điểm hỏi: $P=(100,y,0)$ --- nằm trên mặt phẳng $z=0$.

\paragraph{Bước 1 --- Phần dữ liệu $u_h$ triệt tiêu (đối xứng).}
Phép phản xạ $R:(x,y,z)\mapsto(x,y,-z)$ đổi $B_1(0,0,2)\leftrightarrow B_1(0,0,-2)$, nên
$h(Rp)=-h(p)$ ($h$ \emph{lẻ} qua mặt phẳng $z=0$). Công thức Kirchhoff với $g=0$:
\[
  u_h(P,t)=\frac{1}{4\pi t}\int_{\partial B_t(P)} h\,dS .
\]
Vì $P$ nằm trên $z=0$ nên $RP=P$ và mặt cầu $\partial B_t(P)$ bất biến qua $R$. Đổi biến $q=Rp$:
\[
  \int_{\partial B_t(P)} h\,dS=\int_{\partial B_t(P)} h(Rq)\,dS=-\int_{\partial B_t(P)} h\,dS
  \;\Rightarrow\; \int_{\partial B_t(P)} h\,dS=0 .
\]
\begin{nhobox}
$u_h(100,y,0,t)=0$ với mọi $t$ --- hai quả cầu đối xứng qua mặt phẳng chứa điểm hỏi nên đóng góp khử nhau.
\end{nhobox}

\paragraph{Bước 2 --- Phần nguồn $u_f$ hạ về 1D.}
Nguồn $f=\chi_{[0,1]}(t)\chi_{[-1,1]}(x)$ không phụ thuộc $y,z$. Vì dữ liệu ban đầu của $u_f$
bằng 0 nên $u_f$ cũng không phụ thuộc $y,z$: đặt $u_f(x,y,z,t)=v(x,t)$ thì $u_{f,tt}=v_{tt}$,
$\Delta u_f=v_{xx}$, dẫn về \textbf{phương trình sóng 1D}
\[
  v_{tt}=v_{xx}+\chi_{[0,1]}(t)\chi_{[-1,1]}(x),\quad v(x,0)=0,\ v_t(x,0)=0 .
\]
Công thức d'Alembert (Duhamel 1D):
\[
  v(x,t)=\frac12\int_0^t\!\!\int_{x-(t-s)}^{x+(t-s)} \chi_{[0,1]}(s)\chi_{[-1,1]}(\xi)\,d\xi\,ds .
\]

\paragraph{Bước 3 --- Tính tại $x=100$.}
Đặt $\tau=t-s$. Độ dài giao của $[\,100-\tau,\,100+\tau\,]$ với $[-1,1]$:
\[
  L(\tau)=\begin{cases}
    0, & \tau<99,\\
    \tau-99, & 99\le\tau\le101,\\
    2, & \tau>101.
  \end{cases}
\]
Vì nguồn chỉ bật khi $s\in[0,1]$:
\[
  v(100,t)=\frac12\int_{\max(0,\,t-1)}^{t} L(\tau)\,d\tau
  \stackrel{t\ge1}{=}\frac12\int_{t-1}^{t} L(\tau)\,d\tau .
\]
Tính theo từng đoạn:
\begin{nhobox}
\[
  u(100,y,0,t)=v(100,t)=\begin{cases}
    0, & t\le 99,\\[2pt]
    \dfrac{(t-99)^2}{4}, & 99\le t\le 100,\\[6pt]
    \dfrac{2t-199}{4}, & 100\le t\le 101,\\[6pt]
    (t-100)-\dfrac{(t-100)^2}{4}, & 101\le t\le 102,\\[6pt]
    1, & t\ge 102.
  \end{cases}
\]
Đáp số \emph{không phụ thuộc $y$}; $u_h$ đã khử, chỉ còn phần nguồn 1D. Tín hiệu bắt đầu tại
$t=99$, tăng dần rồi giữ giá trị không đổi $1$ từ $t=102$ (nguồn tắt, dịch chuyển dư trong sóng 1D).
\end{nhobox}

\paragraph{Kiểm tra liên tục.}
$t=100$: $\tfrac{1}{4}$ (cả hai nhánh); $t=101$: $\tfrac34$; $t=102$: $1$. Khớp, $v$ liên tục.

% ---------------------------------------------------------------
\subsection*{Dạng bài lân cận}
% ---------------------------------------------------------------
\begin{description}
  \item[\textbf{[Mạnh]} Cauchy sóng 3D thuần nhất --- Kirchhoff] Bỏ nguồn ($f=0$), chỉ Kirchhoff.
  \item[\textbf{[Mạnh]} Cauchy sóng 3D có nguồn --- Duhamel] Như đề này: Duhamel + Kirchhoff.
  \item[\textbf{[Mạnh]} Cauchy sóng 1D --- d'Alembert] Nguồn/dữ liệu độc lập $y,z$ $\Rightarrow$ hạ về 1D.
  \item[\textbf{[Mạnh]} Ước lượng miền ảnh hưởng tại điểm xa] ``$u$ có bằng 0 tại $x=100$?'' --- hình học nón.
  \item[Sóng dữ liệu gián đoạn (nghiệm yếu)] Dữ liệu không trơn $\Rightarrow$ hiểu theo nghĩa phân bố.
  \item[Cauchy phương trình nhiệt 3D] Parabolic: không Huygens, không vận tốc hữu hạn.
  \item[Laplace/Poisson trong $\mathbb{R}^3$] Elliptic: không thời gian, không lan truyền.
\end{description}
